\section{Topology Preview: Continuity Without Rulers}
\chapterline

\subsection{The origin: from Königsberg bridges to the three-body problem}

In elementary geometry and introductory calculus, every statement leans on a ruler. You measure the distance $|x - y|$ between points, you compute angles between vectors, and you declare a function continuous when small changes in input guarantee small changes in output:
\[
|x - c| < \delta \implies |f(x) - f(c)| < \eps.
\]
Both $\eps$ and $\delta$ are measurements with a ruler.

The earliest hint that geometry could survive without a ruler arrived in 1736, when Leonhard Euler resolved the puzzle of the seven bridges of K\"onigsberg. Euler observed that the exact distances between the bridges, the speed of the river, and the land areas of the islands were completely irrelevant to whether a walker could cross every bridge exactly once. Only the pattern of connectivity---the topology of the network---determined the outcome. Euler called this qualitative inquiry \emph{geometria situs} (the geometry of position).

Yet the crisis that forced the birth of modern topology as a systematic discipline arose over 150 years later, in the work of Henri Poincar\'e. In 1889, King Oscar II of Sweden sponsored a prestigious mathematical competition to answer a profound celestial question: is the solar system stable? Can three gravitational bodies (such as the Sun, Earth, and Moon) orbit each other indefinitely according to Newton's laws, or will one eventually be ejected into deep space or crash into the Sun?

Poincar\'e discovered that the three-body problem is non-integrable: analytical perturbation series diverge, and chaotic orbits produce trajectories of dizzying, unpredictable complexity. Calculating explicit coordinates was impossible. Poincar\'e realized that to determine whether orbits remained bounded, mathematicians did not need coordinates; they needed to understand the \emph{qualitative shape of the phase space}. In his groundbreaking 1895 monograph \emph{Analysis Situs}, Poincar\'e invented the foundational machinery of algebraic topology---homology, Betti numbers, and the fundamental group---to classify multidimensional spaces by their holes and loops.

\subsection{Geometry on a sheet of pliable rubber}

Now imagine a universe made of pliable, stretchable rubber. You can stretch the plane like dough, bend it, compress it, or twist it. All distances are distorted. Angles warp and change. A circle stretches into an ellipse, then into an irregular amoeba shape. A straight line sags into a curve.

What geometric properties survive this kind of distortion?

Distance does not survive. Straightness does not survive. Area does not survive. But something profound remains intact: \emph{nearness}. If two points are connected by an unbroken path, stretching the rubber does not sever that path. If a region has a hole in it, stretching cannot plug the hole without gluing rubber together. And if a point sits safely inside a region with room to wiggle, it still has room to wiggle after the stretch.

Topology is the study of these qualitative geometric properties. It is geometry with the ruler discarded.

\begin{center}
\begin{tikzpicture}[scale=1.1, >=Latex]
  % Region U
  \draw[thick, sourceblue, fill=sourceblue!12] plot [smooth cycle, tension=0.7]
    coordinates {(-3.2, 0) (-2, 1.8) (1.2, 2.0) (2.8, 0.4) (1.8, -1.8) (-1.5, -1.6)};
  \node[font=\bfseries, text=sourceblue] at (-2.0, 1.2) {Region $U$};

  % Interior point p
  \coordinate (P) at (-0.6, 0.3);
  \fill[targetgreen] (P) circle (2pt) node[below=2pt, font=\small, text=targetgreen] {$p$ (interior)};
  \draw[dashed, targetgreen, thick, fill=targetgreen!15] (P) circle (0.75);
  \draw[->, targetgreen] (P) -- ++(45:0.75) node[midway, above left, font=\scriptsize] {$r$};
  \node[above=20pt, font=\scriptsize, text=targetgreen, align=center] at (P) {Cushioned inside $U$:\\$B_r(p) \subset U$};

  % Boundary point q
  \coordinate (Q) at (2.4, 0.9);
  \fill[quotientorange] (Q) circle (2pt) node[above right=2pt, font=\small, text=quotientorange] {$q \in \partial U$ (boundary)};
  \draw[dashed, quotientorange, thick, fill=quotientorange!20] (Q) circle (0.65);
  \node[right=24pt, font=\scriptsize, text=quotientorange, align=left] at (Q) {Every ball around $q$ touches\\both inside and outside};

  % Exterior point z
  \coordinate (Z) at (2.4, -1.5);
  \fill[trapred] (Z) circle (2pt) node[below right=2pt, font=\small, text=trapred] {$z$ (exterior)};
  \draw[dashed, trapred, thick, fill=trapred!15] (Z) circle (0.5);
  \node[below=16pt, font=\scriptsize, text=trapred, align=center] at (Z) {Completely outside $U$};
\end{tikzpicture}
\end{center}

\subsection{The definition of an open set}

To formalize nearness without distances, mathematicians realized that the fundamental building block is the concept of an \textbf{open set}.

In ordinary Euclidean space $\R^n$, an open ball around a center $x$ of radius $r > 0$ is the set of all points strictly closer than $r$:
\[
B_r(x) = \{y \in \R^n : \|y - x\| < r\}.
\]
A subset $U \subseteq \R^n$ is called \textbf{open} if every point inside $U$ has a buffer zone of positive radius completely contained inside $U$:
\[
\forall x \in U, \; \exists r > 0 \text{ such that } B_r(x) \subseteq U.
\]
If you stand at any point $x$ in an open set, you can take a small step in \emph{any} direction without falling off the edge. Points on the boundary do not belong to the set.

\subsection{The topological axioms}

What happens when we discard the metric $\|\cdot\|$ entirely? We keep only the collection $\mathcal{T}$ of subsets that we agree to call ``open''. To behave like genuine neighborhoods, this collection must satisfy three simple rules.

\begin{definition}
A \textbf{topology} on a set $X$ is a collection $\mathcal{T}$ of subsets of $X$ satisfying:
\begin{enumerate}
\item Both the empty set $\emptyset$ and the full space $X$ belong to $\mathcal{T}$.
\item Any union (finite or infinite) of sets in $\mathcal{T}$ belongs to $\mathcal{T}$.
\item Any intersection of \emph{finitely many} sets in $\mathcal{T}$ belongs to $\mathcal{T}$.
\end{enumerate}
The pair $(X, \mathcal{T})$ is called a \textbf{topological space}, and the sets in $\mathcal{T}$ are called the \textbf{open sets} of the space.
\end{definition}

Why do unions permit infinitely many sets, while intersections only permit finitely many?

Consider open intervals in $\R$:
\begin{itemize}
\item If you take the union of any collection of open intervals, every point in the union belongs to at least one interval, where it already has breathing room. That room remains valid in the union.
\item But consider an infinite sequence of shrinking open intervals centered at $0$:
\[
\bigcap_{n=1}^\infty \left(-\frac{1}{n}, \frac{1}{n}\right) = \{0\}.
\]
Each individual interval $(-1/n, 1/n)$ is open. Yet their infinite intersection contains only the solitary point $0$. The set $\{0\}$ has zero wiggle room: any open interval around $0$ immediately leaks into positive and negative numbers. The cushion has been pinched to death by the infinite intersection!
\end{itemize}
Limiting the intersection axiom to finitely many sets preserves the breathing room: if you have finitely many cushions $r_1, r_2, \dots, r_k > 0$, the minimum radius $r = \min(r_1, \dots, r_k)$ is still strictly positive.

\subsection{Continuity reimagined: pulling open sets backward}

In calculus, we say $f$ is continuous if nearby inputs yield nearby outputs. In topology, we translate this without mentioning distance.

\begin{center}
\begin{tikzpicture}[scale=1.0, >=Latex]
  % Space X (Domain)
  \draw[thick, sourceblue, fill=sourceblue!8, rounded corners=12pt] (-4.5, -2) rectangle (-1.2, 2);
  \node[font=\bfseries, text=sourceblue] at (-2.8, 1.6) {Domain $X$};

  % Open set f^{-1}(V) in X
  \draw[thick, targetgreen, fill=targetgreen!20] plot [smooth cycle]
    coordinates {(-3.6, -0.6) (-3.8, 0.8) (-2.0, 1.0) (-1.8, -0.4)};
  \node[font=\small, text=targetgreen, align=center] at (-2.8, 0.2) {Preimage $f^{-1}(V)$\\(\textbf{must be open!})};

  % Forward map arrow f
  \draw[->, very thick, structurepurple] (-1.0, 0.5) -- (1.0, 0.5) node[midway, above=2pt, font=\bfseries] {$f$};

  % Backward pull arrow f^{-1}
  \draw[<-, very thick, quotientorange, dashed] (-1.0, -0.5) -- (1.0, -0.5) node[midway, below=2pt, font=\small] {pull back $f^{-1}$};

  % Space Y (Codomain)
  \draw[thick, sourceblue, fill=sourceblue!8, rounded corners=12pt] (1.2, -2) rectangle (4.5, 2);
  \node[font=\bfseries, text=sourceblue] at (2.8, 1.6) {Codomain $Y$};

  % Open set V in Y
  \draw[thick, quotientorange, fill=quotientorange!20] plot [smooth cycle]
    coordinates {(2.0, -0.8) (1.8, 0.7) (3.6, 0.9) (3.8, -0.5)};
  \node[font=\small, text=quotientorange, align=center] at (2.8, 0.1) {Open set $V \subseteq Y$};
\end{tikzpicture}
\end{center}

\begin{definition}
Let $(X, \mathcal{T}_X)$ and $(Y, \mathcal{T}_Y)$ be topological spaces. A function $f: X \to Y$ is \textbf{continuous} if for every open set $V \in \mathcal{T}_Y$, the inverse image
\[
f^{-1}(V) = \{x \in X : f(x) \in V\}
\]
is an open set in $X$.
\end{definition}

Why does this definition use the \emph{inverse} image $f^{-1}(V)$ rather than the forward image $f(U)$?

Consider the constant function $f: \R \to \R$ defined by $f(x) = 5$ for all $x$. Intuitively, a constant function is as smooth and continuous as a function can possibly be. Yet if you take the open interval $U = (0, 1)$, its forward image is:
\[
f(U) = \{5\},
\]
which is a single point, hence not open! Pushing open sets forward destroys openness whenever a function compresses space.

Pulling backward, however, behaves flawlessly:
\begin{itemize}
\item If $5 \in V$, then $f^{-1}(V) = \R$, which is open.
\item If $5 \notin V$, then $f^{-1}(V) = \emptyset$, which is open.
\end{itemize}
The inverse image respects unions, intersections, and complements without distorting the logic of neighborhoods.

\begin{intuitionbox}
Topological continuity says: if you choose any target neighborhood $V$ around an output, all inputs that land inside $V$ form a cushioned neighborhood with room to move. There are no sudden leaps or tears.
\end{intuitionbox}

\begin{historybox}
The path to understanding qualitative geometry without rulers was filled with surprising false starts. For over a century, mathematicians believed Euler's formula for polyhedra, $V - E + F = 2$, was an unassailable law of three-dimensional solids. But in the nineteenth century, shapes with tunnels (like picture frames) and solids with hollow cavities began to emerge, yielding counts of zero or four. As Imre Lakatos later chronicled, mathematicians spent decades arguing whether to ban these strange shapes as ``monsters'' or to repair the formula by discovering the Euler characteristic: $\chi = 2 - 2g$, where $g$ counts the number of holes.

Even Henri Poincar\'e, the father of modern topology, stumbled on the subtlety of higher-dimensional spaces. In his initial 1895 \emph{Analysis Situs}, Poincar\'e conjectured that any three-dimensional space with the same loop-counting homology as a three-dimensional sphere must actually be a sphere. Nine years later, he discovered his own counterexample: the Poincar\'e homology sphere, built by gluing opposite faces of a solid dodecahedron with a slight twist. It shared the homology of a sphere, yet an ant walking along its paths could trace loops that could never be shrunk to a point. Realizing that homology was too blunt an instrument, Poincar\'e formulated the celebrated Poincar\'e Conjecture, which stood unresolved until Grigori Perelman proved it a century later.

Early pioneers also assumed that topology must rely on metric distances, measuring how far points lay from one another. In 1914, Felix Hausdorff cut the cord entirely: he showed that nearness does not need a ruler at all, defining topological spaces purely through neighborhoods and sets. Refined into the modern three axioms of open sets by Kazimierz Kuratowski, topology emerged as the pure study of nearness without measurement.
\end{historybox}

\subsection{Qualitative geometry---and its natural boundaries}

Discarding the ruler unlocks deep qualitative truths about space. Topological invariants, like the fundamental group and the Euler characteristic, remain invariant no matter how much dough is kneaded or stretched. If a space has a hole that cannot be pinched shut, no continuous deformation can ever erase it. This allows mathematicians to classify all two-dimensional surfaces completely, and provides the natural language for infinite-dimensional function spaces in quantum mechanics where ordinary Euclidean distances fail to capture how functions converge.

Yet topology's greatest strength is also its sharpest limitation: by throwing away the ruler, it throws away quantitative geometry entirely. A pebble, an asteroid, and a squashed lump of clay are all topologically identical to a smooth sphere. Topology cannot tell you which trajectory is the shortest path between two cities, how sharply a road bends around a mountain, or how gravity curves the fabric of space. Bare topological spaces have no tangent vectors, no derivatives, and no differential equations.

To re-introduce calculus onto curved worlds without sacrificing coordinate independence, Chapter 23 introduces \textbf{Differential Geometry Preview: Calculus on Curved Worlds}, equipping manifolds with smooth overlapping charts, tangent planes, and Riemannian metric tensors. Later, Chapter 37 provides our rigorous \textbf{Second Pass: Topology as the Study of Stable Nearness}.

\newpage
